\documentclass[12pt]{article}
\usepackage[utf8]{inputenc}
\usepackage[T1]{fontenc}
\usepackage{hyperref}
\title{Yandjibara Language Resources for K-12 Curriculum Development}
\author{Compiled by Claude Sonnet 5.5 (claude-sonnet-5-5) for Andr\'as Kornai}
\date{}
\begin{document}
\maketitle

\section{Introduction}
Yandjibara (ISO 639-3 xyb) is an extinct Pama-Nyungan language of Queensland, in the Maric branch. It is listed at about 24\textdegree 45'S, 145\textdegree 39'E, in the Barcoo region, and the source cited is Breen (1990). It must not be confused with Yindjibarndi, a living Pilbara language of Western Australia, which has a published wordlist and dictionary. Almost nothing about Yandjibara turned up in the searches beyond a database entry.

\section{Primary Educational Resources}
None found.

\section{Linguistic and Reference Materials}
\begin{itemize}
\item ASJP database entry for Yandjibara: \url{https://asjp.clld.org/languages/YANDJIBARA}
\end{itemize}

\section{Religious and Community Resources}
None found.

\section{Digital and Audio Resources}
None found. Resources for Yindjibarndi, such as the Wangka Maya materials, concern a different language and are not included here.

\section{Curriculum Development Context}
No curriculum can be based on what was found. A next step would be to check Breen (1990) and the AIATSIS Austlang records, and to contact Aboriginal organisations in the Barcoo region of Queensland. The searches did not locate any such group.
\end{document}